\documentclass[10pt,a4paper]{amsart}
\usepackage[margin=19mm]{geometry}
\usepackage{amsmath,amssymb,mathtools}
\usepackage[scr=rsfs,cal=euler]{mathalfa}
\usepackage{xfrac}
\usepackage[unicode,hidelinks,bookmarks=false]{hyperref}
\newcommand{\ie}{i.e.\ }
\newcommand{\cf}{cf.\ }
\newcommand{\resp}{respectively\ }
\newcommand{\Subsection}{\subsection}
\providecommand{\nobreakdash}{\nobreak\hbox{-}}
\setlength{\emergencystretch}{2em}
\multlinegap0pt
\usepackage{bm}
\usepackage{bbm}
\usepackage{dsfont}
\usepackage{xspace}
\usepackage{cancel}
\usepackage{mathtools}
\usepackage{amsthm}
\makeatletter
\@ifundefined{proposition}{\newtheorem{proposition}{Proposition}}{}
\makeatother
\title[Bilinear rough singular integrals under a fractional geometr]{Bilinear rough singular integrals under a fractional geometric condition}
\author{Binwei Dan, Moyan Qin, Qingying Xue}
\date{Source: arXiv:2606.27141v1 (CC BY 4.0)}
\begin{document}
\maketitle

\noindent\emph{Source and licence.} Bilinear rough singular integrals under a fractional geometric condition. Version arXiv:2606.27141v1 (CC BY 4.0), distributed under the Creative Commons Attribution 4.0 International licence, as recorded in the arXiv licence metadata for this exact version and shown on the abstract page https://arxiv.org/abs/2606.27141v1. Full rights notice: licenses/arxiv-2606.27141-CC-BY-4.0.txt.\par\smallskip

\noindent\textit{Editorial scope.} Counted unit: the Proposition whose source label is \texttt{propforward\_1d\_orlicz} in the article (source0.tex lines 1194--1204, source label \texttt{propforward\_1d\_orlicz}), delivered together with the article's own complete original proof of it (source0.tex lines 1206--1342, one \texttt{proof} environment of 137 lines). No mathematics is written, repaired or re-derived: statement and proof are the article's own text line for line. Required context retained uncounted: none. Every definition, display and label of those lines is retained unchanged. Omitted from this package and not redistributed: 1--1193, 1205--1205, 1343--1538. Nothing inside a retained excerpt is omitted or altered: every retained body is delivered exactly as the article's lines are written, so no omission receipt of any kind accompanies this package. Citations: the retained excerpts carry no citation command at all, so no bibliography entry of the article is transcribed and no reference is redistributed. Reference adaptation: none was needed and none was made; every reference inside the delivered unit resolves inside it, so no unresolved reference or citation is produced. Printed slips kept literal: no printed word, symbol, hypothesis or number of the counted statement or of its proof was altered. Layout and class: the article's own class is \texttt{amsart} with packages including graphics, color, geometry, amssymb, enumerate, amsmath, bbm, bm, eso-pic, graphicx, tikz, cite, esint, hyperref; the wrapper is the lane's pinned \texttt{amsart} editorial wrapper at 10pt on A4, which restates the theorem-like environments the retained text needs and adds the article's own macro definitions that the retained text needs (none), with extra wrapper packages (bm, bbm, dsfont, xspace, cancel, mathtools). The delivered numbering is the unit's own. Assets: no retained span contains a figure environment, a graphics-inclusion command, a PDF-page inclusion, a table or a TikZ picture; no image file is copied into this package, the package declares no asset dependency and no image file is redistributed with it. Licence: version arXiv:2606.27141v1 of this article is distributed under the Creative Commons Attribution 4.0 International licence, as recorded in the arXiv licence metadata for this exact version and shown on the abstract page https://arxiv.org/abs/2606.27141v1. A full rights notice is delivered at licenses/arxiv-2606.27141-CC-BY-4.0.txt. Characters: every retained excerpt is pure ASCII, so the delivered engine, XeLaTeX with its default Latin Modern text font, reports no missing character.
\begin{proposition}\label{propforward_1d_orlicz}
	Let $0<a<1$ and $\alpha>1$. There exists a nonnegative function
	$\Omega \in L^1(\mathbb{S}^1)$ such that
	\[
	\Omega\in \mathcal K_a(\mathbb{S}^1)
	\]
	but
	\[
	\Omega\notin L(\log L)^\alpha(\mathbb{S}^1).
	\]
\end{proposition}

\begin{proof}
	We identify $\mathbb{S}^1$ with $[0,2\pi)$ endowed with periodic arc-length measure. For each integer $k \ge 1$, we first partition $\mathbb{S}^1$ into $N_k$ consecutive arcs of equal length $2\pi/N_k$, where the number of partition intervals is given by
	$$N_k = \Bigl\lceil k^{-\alpha}4^{k(\frac{1}{1-a} - 1)}\,k^{\frac{2}{1-a}}\Bigr\rceil.$$
	Due to the dominant exponential growth of $4^{k(\frac{1}{1-a} - 1)}$ (since $0<a<1$), we have $N_k \to \infty$ as $k \to \infty$. The ceiling function naturally guarantees that $N_k \ge 1$. 
	
	In the center of each such arc, we choose a closed subinterval $I_{k,j}$ for $1 \le j \le N_k$. 
	
	Next, we define the height and total mass parameters at level $k$ as
	$$h_k = 4^k, \qquad w_k = k^{-\alpha}4^{-k}.$$
	To uniformly distribute the mass $w_k$ around the circle, we set the length of each individual subinterval to be
	$$\delta_k = \frac{w_k}{N_k},$$
	so that $|I_{k,j}| = \delta_k$. Since $w_k \le 1$ for all $k \ge 1$, it is easy to observe that
	$$\delta_k = \frac{w_k}{N_k} \le \frac{1}{N_k} < \frac{2\pi}{N_k}.$$
	This strict inequality guarantees that the subintervals $I_{k,j}$ are strictly contained within their respective ambient arcs and are therefore pairwise disjoint within the fixed level $k$.
	
	We then define the set $E_k$ as the union of these uniformly distributed intervals
	$$E_k = \bigcup_{j=1}^{N_k} I_{k,j}.$$
	By construction, the measure of this set is precisely $\sigma(E_k) = N_k \delta_k = w_k$.
	
	Finally, having constructed the sets $E_k$ for all $k \ge 1$, we define the non-negative measurable function $\Omega$ on $\mathbb{S}^1$ by
	$$\Omega(\theta) = \sum_{k=1}^\infty h_k \chi_{E_k}(\theta).$$
	We now verify that this $\Omega$ satisfies all the required properties.
	
	\medskip
	\noindent\textbf{Step 1 \(\Omega\in L^1(\mathbb{S}^1)\).}
	
	Since $\Omega\ge 0$, Tonelli's theorem gives
	\[
	\int_{\mathbb{S}^1}\Omega(\theta)\,d\sigma(\theta)
	=\sum_{k=1}^\infty h_k \sigma(E_k)
	=\sum_{k=1}^\infty h_k w_k
	=\sum_{k=1}^\infty 4^k\cdot k^{-\alpha}4^{-k}
	=\sum_{k=1}^\infty k^{-\alpha}<\infty,
	\]
	because $\alpha>1$. Hence $\Omega\in L^1(\mathbb{S}^1)$.
	
	\medskip
	\noindent\textbf{Step 2 \(\Omega\notin L(\log L)^\alpha(\mathbb{S}^1)\).}
	
	Define
	\[
	A_k=E_k\setminus \bigcup_{m>k}E_m.
	\]
	Then
	\[
	\sigma(A_k)\ge \sigma(E_k)-\sum_{m>k}\sigma(E_m)
	=w_k-\sum_{m>k}w_m.
	\]
	Since $w_m=m^{-\alpha}4^{-m}$ and $m\mapsto m^{-\alpha}$ is decreasing,
	\[
	\sum_{m>k}w_m
	\le \frac{1}{(k+1)^\alpha}\sum_{m=k+1}^\infty 4^{-m}
	=\frac{1}{(k+1)^\alpha}\frac{4^{-k-1}}{1-1/4}
	=\frac{1}{3}(k+1)^{-\alpha}4^{-k}
	\le \frac{1}{3}k^{-\alpha}4^{-k}
	=\frac{1}{3}w_k.
	\]
	Therefore
	\[
	\sigma(A_k)\ge \frac{2}{3}w_k.
	\]
	If $\theta\in A_k$, then $\theta\in E_k$ and $\theta\notin E_m$ for all $m>k$, hence
	\[
	\Omega(\theta)=\sum_{m=1}^\infty h_m\chi_{E_m}(\theta)\ge h_k.
	\]
	It follows that
	\begin{align*}
		\int_{\mathbb{S}^1}\Omega(\theta)\bigl(\log(e+\Omega(\theta))\bigr)^\alpha\,d\sigma(\theta)
		&\ge \sum_{k=1}^\infty
		\int_{A_k} h_k\bigl(\log(e+h_k)\bigr)^\alpha\,d\sigma(\theta) \\
		&= \sum_{k=1}^\infty h_k\bigl(\log(e+h_k)\bigr)^\alpha \sigma(A_k).
	\end{align*}
	Using $h_k=4^k$, $\sigma(A_k)\ge \frac23 w_k$, and
	$\log(e+4^k)\ge k\log 4$, we obtain
	\begin{align*}
		\int_{\mathbb{S}^1}\Omega(\theta)\bigl(\log(e+\Omega(\theta))\bigr)^\alpha\,d\sigma(\theta)
		&\ge \frac{2}{3}\sum_{k=1}^\infty
		4^k (k\log 4)^\alpha w_k \\
		&= \frac{2}{3}(\log 4)^\alpha
		\sum_{k=1}^\infty 4^k k^\alpha \cdot k^{-\alpha}4^{-k} \\
		&= \frac{2}{3}(\log 4)^\alpha \sum_{k=1}^\infty 1
		= \infty.
	\end{align*}
	Therefore $\Omega\notin L(\log L)^\alpha(\mathbb{S}^1)$.
	
	\medskip
	\noindent\textbf{Step 3 \(\Omega\in \mathcal K_a(\mathbb{S}^1)\).}
	
	For each $k\ge 1$, define
	$$ I_k(\xi) = \int_{E_k} \frac{1}{|\theta \cdot \xi|^a} \, d\sigma(\theta). $$
	Then, by Tonelli's theorem,
	$$ \int_{\mathbb{S}^1} \frac{\Omega(\theta)}{|\theta \cdot \xi|^a} \, d\sigma(\theta) = \sum_{k=1}^\infty h_k I_k(\xi). $$
	Hence it suffices to prove that
	$$ \sup_{\xi\in\mathbb{S}^1} I_k(\xi) \lesssim_a \delta_k^{\,1-a} + w_k \qquad \text{uniformly in } k. $$
	
	Fix $\xi \in \mathbb{S}^1$ and define its orthogonal poles $P_\xi = \{\xi^\perp, -\xi^\perp\}$. For any $\theta \in \mathbb{S}^1$, the inner product satisfies
	$$ |\theta \cdot \xi| = \sin(d_{\mathbb{S}^1}(\theta, P_\xi)), \quad \text{where } d_{\mathbb{S}^1}(\theta, P_\xi) = \min_{p \in P_\xi} \arccos(\theta \cdot p). $$
	
	Recall that $\mathbb{S}^1$ is partitioned into $N_k$ arcs, with each subinterval $I_{k,j}$ centered in its ambient arc. We define the set of \emph{bad} intervals as
	$$ \mathcal{B}_k(\xi) = \{ I_{k,j}:I_{k,j} \cap P_\xi \neq \emptyset \}, $$
	and call the remaining intervals \emph{good}. Since there are exactly two poles in $P_\xi$, and each pole can intersect at most two adjacent partition arcs, there are at most $4$ bad intervals in $\mathcal{B}_k(\xi)$.
	
	For any bad interval $I_{k,j} \in \mathcal{B}_k(\xi)$, its ambient arc intersects a pole $p \in P_\xi$, introducing a singularity. Parametrizing $\theta \in I_{k,j}$ by its distance to the nearest pole, $\varphi = d_{\mathbb{S}^1}(\theta, P_\xi)$, gives $|\theta \cdot \xi| = \sin \varphi$. This maps $I_{k,j}$ to an angular interval $\widetilde{I}_{k,j}$. On the unit circle, the arc length equals the magnitude of its corresponding central angle, so $|\widetilde{I}_{k,j}| = |I_{k,j}| = \delta_k$ and $d\sigma(\theta) = d\varphi$. Thus, changing variables yields
	$$ \int_{I_{k,j}} |\theta \cdot \xi|^{-a} \, d\sigma(\theta) = \int_{\widetilde{I}_{k,j}} |\sin \varphi|^{-a} \, d\varphi. $$
	
	Observe that $p \in I_{k,j}$ is equivalent to $0 \in \widetilde{I}_{k,j}$ in the parameter domain. Because the value of $|\sin \varphi|^{-a}$ gets larger as it gets closer to $0$, the integral $\int_{\widetilde{I}_{k,j}} |\sin \varphi|^{-a} \, d\varphi$ is maximized when $0$ is positioned exactly at the center of $\widetilde{I}_{k,j}$. Therefore, we have
	$$\int_{\widetilde{I}_{k,j}} |\sin \varphi|^{-a} \, d\varphi \le \int_{-\delta_k/2}^{\delta_k/2} |\sin \varphi|^{-a} \, d\varphi.$$
	
	Since $\delta_k/2 \le \pi/2$, we have $|\sin \varphi| \ge \frac{2}{\pi} |\varphi|$, which yields
	$$ \int_{-\delta_k/2}^{\delta_k/2} |\sin \varphi|^{-a} \, d\varphi \le 2 \left(\frac{\pi}{2}\right)^a \int_0^{\delta_k/2} \varphi^{-a} \, d\varphi \lesssim_a \delta_k^{\,1-a}. $$
	Summing over the bounded number of bad intervals gives
	$$ \sum_{I_{k,j} \, \text{bad}} \int_{I_{k,j}} |\theta \cdot \xi|^{-a} \, d\sigma(\theta) \lesssim_a \delta_k^{\,1-a}.$$
	
	We call the remaining intervals \emph{good}. For $1 \le m \le \lfloor N_k/4 \rfloor$, we group these good intervals by their distance to the poles $P_\xi$ by defining
	$$ \mathcal{G}_{k,m}(\xi) = \left\{ I_{k,j} \notin \text{bad}  \frac{2\pi m}{N_k} \le d_{\mathbb{S}^1}(I_{k,j}, P_\xi) < \frac{2\pi (m+1)}{N_k} \right\}. $$
	By construction, each $\mathcal{G}_{k,m}$ contains at most 4 intervals $I_{k,j}$, as each of the two poles admits at most one corresponding interval in both the clockwise and counterclockwise directions at a given grid distance $m$. Furthermore, the distance from any point in these intervals to the nearest singularity is strictly bounded from below by
	$$ \inf_{\theta \in I_{k,j} \in \mathcal{G}_{k,m}(\xi)} d_{\mathbb{S}^1}(\theta, P_\xi) \ge \frac{2\pi m}{N_k} \gtrsim \frac{m}{N_k}. $$
	
	Consequently, for any interval $I_{k,j} \in \mathcal{G}_{k,m}(\xi)$, we can uniformly bound the singular kernel by
	$$ |\theta \cdot \xi|^{-a} = |\sin(d_{\mathbb{S}^1}(\theta, P_\xi))|^{-a} \lesssim \left(\frac{N_k}{m}\right)^a.$$
	Integrating this bound over the length $\delta_k$ of the interval gives
	$$ \int_{I_{k,j}} |\theta \cdot \xi|^{-a} \, d\sigma(\theta) \lesssim \delta_k \left(\frac{N_k}{m}\right)^a. $$
	Summing over the at most 4 intervals in $\mathcal{G}_{k,m}(\xi)$ and then over all possible discrete distances $m$ yields
	$$ \sum_{I_{k,j} \, \text{good}} \int_{I_{k,j}} |\theta \cdot \xi|^{-a} \, d\sigma(\theta) \lesssim \delta_k N_k^a \sum_{m=1}^{\lfloor N_k/4 \rfloor} m^{-a}. $$
	Since $0 < a < 1$, the sum is asymptotically comparable to $N_k^{1-a}$, and since $\delta_k = w_k/N_k$, we have
	$$ \sum_{I_{k,j} \, \text{good}} \int_{I_{k,j}} |\theta \cdot \xi|^{-a} \, d\sigma(\theta) \lesssim_a \delta_k N_k^a \cdot N_k^{1-a} = \delta_k N_k = w_k. $$
	
	Summing the contributions from both the bad and good intervals gives
	$$ I_k(\xi) = \int_{E_k} \frac{1}{|\theta \cdot \xi|^a} \, d\sigma(\theta) \lesssim_a \delta_k^{1-a} + w_k. $$
	Crucially, this uniform bound holds independently of the choice of $\xi \in \mathbb{S}^1$.
	
	Finally, taking the supremum over $\xi$ and summing over all $k$, we conclude
	$$ \sup_{\xi \in \mathbb{S}^1} \int_{\mathbb{S}^1} \frac{|\Omega(\theta)|}{|\theta \cdot \xi|^a} \, d\sigma(\theta) \lesssim_a \sum_{k=1}^\infty h_k \delta_k^{1-a} + \sum_{k=1}^\infty h_k w_k. $$
	The second sum converges because $\sum_{k=1}^\infty h_k w_k = \sum_{k=1}^\infty k^{-\alpha} < \infty$. For the first sum, since $N_k \ge w_k 4^{\frac{k}{1-a}} k^{\frac{2}{1-a}}$, we have $\delta_k = \frac{w_k}{N_k} \le 4^{-\frac{k}{1-a}} k^{-\frac{2}{1-a}}$, which implies $\delta_k^{1-a} \le 4^{-k} k^{-2}$. Thus,
	$$ \sum_{k=1}^\infty h_k \delta_k^{1-a} \le \sum_{k=1}^\infty 4^k \cdot 4^{-k} k^{-2} = \sum_{k=1}^\infty k^{-2} < \infty. $$
	Hence, we conclude that $\sup_{\xi \in \mathbb{S}^1} \int_{\mathbb{S}^1} \frac{|\Omega(\theta)|}{|\theta \cdot \xi|^a} \, d\sigma(\theta) < \infty$, completing the proof of Proposition~\ref{propforward_1d_orlicz}.
\end{proof}

\end{document}
